\documentclass[12pt]{article}
\usepackage{amsmath,amssymb}
\usepackage{geometry}
\usepackage{enumitem}
\geometry{margin=0.85in}

\title{Calculus Practice: Area and Volume Methods}
\author{}
\date{}

\begin{document}
\maketitle

\section*{Formula Reference Sheet}

\subsection*{Area Between Two Curves}
Use vertical slices for $dx$ and horizontal slices for $dy$:
\[
 A=\int_a^b\bigl(\text{top }y-\text{bottom }y\bigr)\,dx,
 \qquad
 A=\int_c^d\bigl(\text{right }x-\text{left }x\bigr)\,dy.
\]
Find intersection points first. Split the integral whenever the top, bottom, right, or left boundary changes. Area is nonnegative, so reverse a subtraction or split at a crossing rather than keeping negative signed area.

\subsection*{Washers and Disks}
A slice perpendicular to the axis of rotation produces a washer or disk:
\[
 V=\pi\int\bigl(R^2-r^2\bigr)\,d(\text{variable}).
\]
Here $R$ is the distance from the axis to the farther boundary and $r$ is the distance to the nearer boundary. For a disk, $r=0$.

\[
\begin{array}{c|c|c}
\text{Axis of rotation} & \text{Washer slice} & \text{Usual variable}\\ \hline
\text{horizontal} & \text{vertical} & dx\\
\text{vertical} & \text{horizontal} & dy
\end{array}
\]
For rotation about $y=k$, radii are vertical distances such as $|f(x)-k|$. For rotation about $x=h$, radii are horizontal distances such as $|g(y)-h|$.

\subsection*{Cylindrical Shells}
A slice parallel to the axis of rotation produces a shell:
\[
 V=2\pi\int(\text{radius})(\text{height})\,d(\text{variable}).
\]
The radius is the distance from the slice to the axis. Shell height is top minus bottom for vertical shells and right minus left for horizontal shells.

\[
\begin{array}{c|c|c}
\text{Axis of rotation} & \text{Shell slice} & \text{Usual variable}\\ \hline
\text{vertical} & \text{vertical} & dx\\
\text{horizontal} & \text{horizontal} & dy
\end{array}
\]

\subsection*{Known Cross Sections}
If every slice has a specified shape, first write its area $A$ in terms of the base width, then integrate:
\[
 V=\int A(x)\,dx
 \qquad\text{or}\qquad
 V=\int A(y)\,dy.
\]
Common cross-sectional areas are
\[
\begin{array}{c|c}
\text{Shape and given measurement} & \text{Area}\\ \hline
\text{Square, side }s & s^2\\
\text{Semicircle, diameter }d & \dfrac{\pi d^2}{8}\\
\text{Equilateral triangle, side }s & \dfrac{\sqrt3}{4}s^2\\
\text{Isosceles right triangle, leg }\ell & \dfrac12\ell^2\\
\text{Isosceles right triangle, hypotenuse }h & \dfrac{h^2}{4}
\end{array}
\]
The base width is found by top minus bottom for vertical slices or right minus left for horizontal slices.

\subsection*{Choosing a Method}
\begin{enumerate}[leftmargin=1.6em,itemsep=0.2em]
  \item Sketch the region and axis of rotation.
  \item Decide whether vertical or horizontal slices avoid splitting or inverse functions.
  \item Perpendicular to the axis means washers; parallel to the axis means shells.
  \item Write every radius, height, or cross-sectional measurement as a distance.
  \item Check bounds and units before evaluating.
\end{enumerate}

\clearpage
\section*{Instructions}
Complete in 60 minutes. Sketch each region, show the integral setup, and evaluate exactly. Choose $dx$ or $dy$ deliberately. For rotation problems, identify whether your slices produce washers/disks or shells before integrating.

\section*{Problems}
\begin{enumerate}[leftmargin=*,itemsep=1.1em]
  \item Find the area enclosed by $y=2x$ and $y=x^2$. Integrate with respect to $x$.

  \item Find the area enclosed by $x=y^2$ and $x=2y$. Integrate with respect to $y$.

  \item Find the total area enclosed by $y=x^3$ and $y=x$. Your setup must account for the change in the upper curve.

  \item The region bounded by $y=x^2$ and $y=2x$ is revolved about the $x$-axis. Find the volume using washers.

  \item The same region bounded by $y=x^2$ and $y=2x$ is revolved about the $y$-axis. Find the volume using shells.

  \item The region bounded by $x=y^2$ and $x=2y$ is revolved about the $y$-axis. Find the volume using washers and integration with respect to $y$.

  \item The region bounded by $y=\sqrt{x}$, $y=0$, and $x=4$ is revolved about the vertical line $x=5$. Find the volume using shells.

  \item In the first quadrant, the region bounded by $y=x^2$ and $x=y^2$ is revolved about the line $y=-1$. Find the volume using washers.

  \item The base of a solid is the region bounded by $y=x$ and $y=x^2$. Cross sections perpendicular to the $x$-axis are squares. Find the volume.

  \item The base of a solid is the region bounded by $y=x$ and $y=x^2$. Cross sections perpendicular to the $x$-axis are semicircles whose diameters lie in the base. Find the volume.

  \item The base of a solid is the region bounded by $x=y^2$ and $x=4$. Cross sections perpendicular to the $y$-axis are equilateral triangles. Find the volume.

  \item The base of a solid is the region bounded by $y=4-x^2$ and $y=0$. Cross sections perpendicular to the $x$-axis are isosceles right triangles whose legs lie in the base. Find the volume.
\end{enumerate}

\end{document}
